%% ══════════════════════════════════════════════════════════════════
%%  astro-nova / references / preamble.tex
%%  NovaForge 科研模式导言区 — 用于论文笔记、文献整理
%%  用法：%% ══════════════════════════════════════════════════════════════════
%%  astro-nova / references / preamble.tex
%%  NovaForge 科研模式导言区 — 用于论文笔记、文献整理
%%  用法：%% ══════════════════════════════════════════════════════════════════
%%  astro-nova / references / preamble.tex
%%  NovaForge 科研模式导言区 — 用于论文笔记、文献整理
%%  用法：%% ══════════════════════════════════════════════════════════════════
%%  astro-nova / references / preamble.tex
%%  NovaForge 科研模式导言区 — 用于论文笔记、文献整理
%%  用法：\input{preamble.tex}
%% ══════════════════════════════════════════════════════════════════

%% ── 中文 ──
\usepackage{xeCJK}
\setCJKmainfont{SimSun}
\usepackage{indentfirst}

%% ── 数学/公式 ──
\usepackage{amsmath,amssymb,amsthm}
\usepackage{cancel}
\usepackage{physics}
\usepackage{mathrsfs}
\usepackage{bm}

%% ── 页面 ──
\usepackage[top=1.4cm,bottom=1.2cm,left=1.0cm,right=1.0cm,includehead]{geometry}
\usepackage{xcolor}
\usepackage{enumitem}
\usepackage{array,booktabs,multirow}
\usepackage{fancyhdr}
\usepackage{titlesec}
\usepackage{hyperref}
\usepackage{environ}
\usepackage{tikz}
\usetikzlibrary{decorations.pathmorphing,patterns,arrows.meta}
\hypersetup{colorlinks=true,linkcolor=blue!60!black,urlcolor=blue!60!black}

%% ── 颜色系统 ──
\definecolor{titlecolor}{HTML}{1a237e}
\definecolor{seccolor}{HTML}{004d40}
\definecolor{subseccolor}{HTML}{b71c1c}
\definecolor{emphcolor}{HTML}{bf360c}
\definecolor{infocolor}{HTML}{1565c0}
\definecolor{warncolor}{HTML}{e65100}
\definecolor{supercolor}{HTML}{880e4f}
\definecolor{examplecolor}{HTML}{1b5e20}
\definecolor{sectionbg}{HTML}{e8eaf6}
\definecolor{lightgray}{HTML}{f5f5f5}
\definecolor{hwcolor}{HTML}{795548}
\definecolor{learncolor}{HTML}{2e7d32}
\definecolor{papercolor}{HTML}{4a148c}   % 文献卡片紫

%% ── 标题格式 ──
\titleformat{\section}{\LARGE\bfseries\color{titlecolor}}{}{0em}{}[\vspace{-0.2em}\rule{\textwidth}{0.8pt}]
\titleformat{\subsection}{\large\bfseries\color{seccolor}}{}{0em}{}
\titleformat{\subsubsection}{\normalsize\bfseries\color{subseccolor}}{}{0em}{}
\titlespacing*{\section}{0pt}{1em}{0.5em}
\titlespacing*{\subsection}{0pt}{0.6em}{0.2em}
\titlespacing*{\subsubsection}{0pt}{0.4em}{0.1em}

%% ── 页眉（使用时替换标题） ──
\pagestyle{fancy}
\fancyhf{}
\fancyhead[L]{\small\color{gray}astro-nova · 文献笔记}
\fancyhead[R]{\small\color{gray}\today}
\fancyfoot[C]{\small\color{gray}\thepage}
\renewcommand{\headrulewidth}{0.4pt}
\setlength{\headheight}{14pt}

%% ── NovaForge 通用命令 ──
\newcommand{\key}[1]{\textcolor{emphcolor}{\textbf{#1}}}
\newcommand{\super}[1]{\textcolor{supercolor}{\textbf{#1}}}
\newcommand{\source}[1]{\textcolor{purple!60!black}{\small [#1]}}
\newcommand{\formula}[1]{\vspace{0.3em}\begin{center}\fcolorbox{titlecolor}{white}{\parbox{0.92\textwidth}{\centering\color{titlecolor}\small #1}}\end{center}\vspace{0.2em}}
\newcommand{\infobox}[1]{\vspace{0.3em}{\small\color{infocolor}\noindent$\blacktriangleright$ #1}\vspace{0.2em}}
\newcommand{\warning}[1]{\vspace{0.3em}{\small\color{warncolor}\noindent\textbf{!} #1}\vspace{0.2em}}
\newcommand{\sep}{\vspace{0.5em}\hrule\vspace{0.5em}}

\newcommand{\knowtitle}[1]{%
  \vspace{0.4em}%
  \noindent\colorbox{sectionbg}{\parbox{\dimexpr\textwidth-2\fboxsep\relax}{\small\bfseries\color{titlecolor}#1}}%
  \vspace{0.2em}%
}

%% ── 科研模式专有命令 ──

% 文献卡片
\newcommand{\paperinfo}[4]{%
  \vspace{0.5em}%
  \noindent
  \begin{tabular}{|p{2.2em}|p{\dimexpr\textwidth-2.2em-2\tabcolsep\relax}|}
    \hline
    \multicolumn{1}{|c|}{\textbf{\small\color{papercolor}标题}} & {\small #1}\\
    \hline
    \multicolumn{1}{|c|}{\textbf{\small\color{papercolor}作者}} & {\small #2}\\
    \hline
    \multicolumn{1}{|c|}{\textbf{\small\color{papercolor}刊源}} & {\small #3}\\
    \hline
    \multicolumn{1}{|c|}{\textbf{\small\color{papercolor}年份}} & {\small #4}\\
    \hline
  \end{tabular}\par
  \vspace{0.3em}%
}

% 科研小标题
\newcommand{\lithead}[1]{%
  \vspace{0.3em}%
  \noindent{\small\bfseries\color{titlecolor}#1}\par
  \vspace{0.1em}%
}

%% ── 环境 ──
\newenvironment{exampenv}[1]{%
  \vspace{0.5em}%
  \noindent{\color{supercolor!15}\rule{0.4em}{0pt}}\hspace{-0.4em}%
  \textcolor{supercolor}{\rule{0.4em}{0.9em}}\hspace{0.3em}%
  \textcolor{supercolor}{\textbf{【案例/例题】#1}}\par
  \noindent\hspace{0.7em}{\small\color{infocolor}分析：}\normalsize
}{%
  \par\vspace{0.3em}%
}

\newenvironment{pracenv}[1]{%
  \vspace{0.5em}%
  \noindent{\color{learncolor!15}\rule{0.4em}{0pt}}\hspace{-0.4em}%
  \textcolor{learncolor}{\rule{0.4em}{0.9em}}\hspace{0.3em}%
  \textcolor{learncolor}{\textbf{【练习/复盘】#1}}\par
  \noindent\hspace{0.7em}{\small\color{gray}解答：}\normalsize
}{%
  \par\vspace{0.3em}%
}

%% ══════════════════════════════════════════════════════════════════

%% ── 中文 ──
\usepackage{xeCJK}
\setCJKmainfont{SimSun}
\usepackage{indentfirst}

%% ── 数学/公式 ──
\usepackage{amsmath,amssymb,amsthm}
\usepackage{cancel}
\usepackage{physics}
\usepackage{mathrsfs}
\usepackage{bm}

%% ── 页面 ──
\usepackage[top=1.4cm,bottom=1.2cm,left=1.0cm,right=1.0cm,includehead]{geometry}
\usepackage{xcolor}
\usepackage{enumitem}
\usepackage{array,booktabs,multirow}
\usepackage{fancyhdr}
\usepackage{titlesec}
\usepackage{hyperref}
\usepackage{environ}
\usepackage{tikz}
\usetikzlibrary{decorations.pathmorphing,patterns,arrows.meta}
\hypersetup{colorlinks=true,linkcolor=blue!60!black,urlcolor=blue!60!black}

%% ── 颜色系统 ──
\definecolor{titlecolor}{HTML}{1a237e}
\definecolor{seccolor}{HTML}{004d40}
\definecolor{subseccolor}{HTML}{b71c1c}
\definecolor{emphcolor}{HTML}{bf360c}
\definecolor{infocolor}{HTML}{1565c0}
\definecolor{warncolor}{HTML}{e65100}
\definecolor{supercolor}{HTML}{880e4f}
\definecolor{examplecolor}{HTML}{1b5e20}
\definecolor{sectionbg}{HTML}{e8eaf6}
\definecolor{lightgray}{HTML}{f5f5f5}
\definecolor{hwcolor}{HTML}{795548}
\definecolor{learncolor}{HTML}{2e7d32}
\definecolor{papercolor}{HTML}{4a148c}   % 文献卡片紫

%% ── 标题格式 ──
\titleformat{\section}{\LARGE\bfseries\color{titlecolor}}{}{0em}{}[\vspace{-0.2em}\rule{\textwidth}{0.8pt}]
\titleformat{\subsection}{\large\bfseries\color{seccolor}}{}{0em}{}
\titleformat{\subsubsection}{\normalsize\bfseries\color{subseccolor}}{}{0em}{}
\titlespacing*{\section}{0pt}{1em}{0.5em}
\titlespacing*{\subsection}{0pt}{0.6em}{0.2em}
\titlespacing*{\subsubsection}{0pt}{0.4em}{0.1em}

%% ── 页眉（使用时替换标题） ──
\pagestyle{fancy}
\fancyhf{}
\fancyhead[L]{\small\color{gray}astro-nova · 文献笔记}
\fancyhead[R]{\small\color{gray}\today}
\fancyfoot[C]{\small\color{gray}\thepage}
\renewcommand{\headrulewidth}{0.4pt}
\setlength{\headheight}{14pt}

%% ── NovaForge 通用命令 ──
\newcommand{\key}[1]{\textcolor{emphcolor}{\textbf{#1}}}
\newcommand{\super}[1]{\textcolor{supercolor}{\textbf{#1}}}
\newcommand{\source}[1]{\textcolor{purple!60!black}{\small [#1]}}
\newcommand{\formula}[1]{\vspace{0.3em}\begin{center}\fcolorbox{titlecolor}{white}{\parbox{0.92\textwidth}{\centering\color{titlecolor}\small #1}}\end{center}\vspace{0.2em}}
\newcommand{\infobox}[1]{\vspace{0.3em}{\small\color{infocolor}\noindent$\blacktriangleright$ #1}\vspace{0.2em}}
\newcommand{\warning}[1]{\vspace{0.3em}{\small\color{warncolor}\noindent\textbf{!} #1}\vspace{0.2em}}
\newcommand{\sep}{\vspace{0.5em}\hrule\vspace{0.5em}}

\newcommand{\knowtitle}[1]{%
  \vspace{0.4em}%
  \noindent\colorbox{sectionbg}{\parbox{\dimexpr\textwidth-2\fboxsep\relax}{\small\bfseries\color{titlecolor}#1}}%
  \vspace{0.2em}%
}

%% ── 科研模式专有命令 ──

% 文献卡片
\newcommand{\paperinfo}[4]{%
  \vspace{0.5em}%
  \noindent
  \begin{tabular}{|p{2.2em}|p{\dimexpr\textwidth-2.2em-2\tabcolsep\relax}|}
    \hline
    \multicolumn{1}{|c|}{\textbf{\small\color{papercolor}标题}} & {\small #1}\\
    \hline
    \multicolumn{1}{|c|}{\textbf{\small\color{papercolor}作者}} & {\small #2}\\
    \hline
    \multicolumn{1}{|c|}{\textbf{\small\color{papercolor}刊源}} & {\small #3}\\
    \hline
    \multicolumn{1}{|c|}{\textbf{\small\color{papercolor}年份}} & {\small #4}\\
    \hline
  \end{tabular}\par
  \vspace{0.3em}%
}

% 科研小标题
\newcommand{\lithead}[1]{%
  \vspace{0.3em}%
  \noindent{\small\bfseries\color{titlecolor}#1}\par
  \vspace{0.1em}%
}

%% ── 环境 ──
\newenvironment{exampenv}[1]{%
  \vspace{0.5em}%
  \noindent{\color{supercolor!15}\rule{0.4em}{0pt}}\hspace{-0.4em}%
  \textcolor{supercolor}{\rule{0.4em}{0.9em}}\hspace{0.3em}%
  \textcolor{supercolor}{\textbf{【案例/例题】#1}}\par
  \noindent\hspace{0.7em}{\small\color{infocolor}分析：}\normalsize
}{%
  \par\vspace{0.3em}%
}

\newenvironment{pracenv}[1]{%
  \vspace{0.5em}%
  \noindent{\color{learncolor!15}\rule{0.4em}{0pt}}\hspace{-0.4em}%
  \textcolor{learncolor}{\rule{0.4em}{0.9em}}\hspace{0.3em}%
  \textcolor{learncolor}{\textbf{【练习/复盘】#1}}\par
  \noindent\hspace{0.7em}{\small\color{gray}解答：}\normalsize
}{%
  \par\vspace{0.3em}%
}

%% ══════════════════════════════════════════════════════════════════

%% ── 中文 ──
\usepackage{xeCJK}
\setCJKmainfont{SimSun}
\usepackage{indentfirst}

%% ── 数学/公式 ──
\usepackage{amsmath,amssymb,amsthm}
\usepackage{cancel}
\usepackage{physics}
\usepackage{mathrsfs}
\usepackage{bm}

%% ── 页面 ──
\usepackage[top=1.4cm,bottom=1.2cm,left=1.0cm,right=1.0cm,includehead]{geometry}
\usepackage{xcolor}
\usepackage{enumitem}
\usepackage{array,booktabs,multirow}
\usepackage{fancyhdr}
\usepackage{titlesec}
\usepackage{hyperref}
\usepackage{environ}
\usepackage{tikz}
\usetikzlibrary{decorations.pathmorphing,patterns,arrows.meta}
\hypersetup{colorlinks=true,linkcolor=blue!60!black,urlcolor=blue!60!black}

%% ── 颜色系统 ──
\definecolor{titlecolor}{HTML}{1a237e}
\definecolor{seccolor}{HTML}{004d40}
\definecolor{subseccolor}{HTML}{b71c1c}
\definecolor{emphcolor}{HTML}{bf360c}
\definecolor{infocolor}{HTML}{1565c0}
\definecolor{warncolor}{HTML}{e65100}
\definecolor{supercolor}{HTML}{880e4f}
\definecolor{examplecolor}{HTML}{1b5e20}
\definecolor{sectionbg}{HTML}{e8eaf6}
\definecolor{lightgray}{HTML}{f5f5f5}
\definecolor{hwcolor}{HTML}{795548}
\definecolor{learncolor}{HTML}{2e7d32}
\definecolor{papercolor}{HTML}{4a148c}   % 文献卡片紫

%% ── 标题格式 ──
\titleformat{\section}{\LARGE\bfseries\color{titlecolor}}{}{0em}{}[\vspace{-0.2em}\rule{\textwidth}{0.8pt}]
\titleformat{\subsection}{\large\bfseries\color{seccolor}}{}{0em}{}
\titleformat{\subsubsection}{\normalsize\bfseries\color{subseccolor}}{}{0em}{}
\titlespacing*{\section}{0pt}{1em}{0.5em}
\titlespacing*{\subsection}{0pt}{0.6em}{0.2em}
\titlespacing*{\subsubsection}{0pt}{0.4em}{0.1em}

%% ── 页眉（使用时替换标题） ──
\pagestyle{fancy}
\fancyhf{}
\fancyhead[L]{\small\color{gray}astro-nova · 文献笔记}
\fancyhead[R]{\small\color{gray}\today}
\fancyfoot[C]{\small\color{gray}\thepage}
\renewcommand{\headrulewidth}{0.4pt}
\setlength{\headheight}{14pt}

%% ── NovaForge 通用命令 ──
\newcommand{\key}[1]{\textcolor{emphcolor}{\textbf{#1}}}
\newcommand{\super}[1]{\textcolor{supercolor}{\textbf{#1}}}
\newcommand{\source}[1]{\textcolor{purple!60!black}{\small [#1]}}
\newcommand{\formula}[1]{\vspace{0.3em}\begin{center}\fcolorbox{titlecolor}{white}{\parbox{0.92\textwidth}{\centering\color{titlecolor}\small #1}}\end{center}\vspace{0.2em}}
\newcommand{\infobox}[1]{\vspace{0.3em}{\small\color{infocolor}\noindent$\blacktriangleright$ #1}\vspace{0.2em}}
\newcommand{\warning}[1]{\vspace{0.3em}{\small\color{warncolor}\noindent\textbf{!} #1}\vspace{0.2em}}
\newcommand{\sep}{\vspace{0.5em}\hrule\vspace{0.5em}}

\newcommand{\knowtitle}[1]{%
  \vspace{0.4em}%
  \noindent\colorbox{sectionbg}{\parbox{\dimexpr\textwidth-2\fboxsep\relax}{\small\bfseries\color{titlecolor}#1}}%
  \vspace{0.2em}%
}

%% ── 科研模式专有命令 ──

% 文献卡片
\newcommand{\paperinfo}[4]{%
  \vspace{0.5em}%
  \noindent
  \begin{tabular}{|p{2.2em}|p{\dimexpr\textwidth-2.2em-2\tabcolsep\relax}|}
    \hline
    \multicolumn{1}{|c|}{\textbf{\small\color{papercolor}标题}} & {\small #1}\\
    \hline
    \multicolumn{1}{|c|}{\textbf{\small\color{papercolor}作者}} & {\small #2}\\
    \hline
    \multicolumn{1}{|c|}{\textbf{\small\color{papercolor}刊源}} & {\small #3}\\
    \hline
    \multicolumn{1}{|c|}{\textbf{\small\color{papercolor}年份}} & {\small #4}\\
    \hline
  \end{tabular}\par
  \vspace{0.3em}%
}

% 科研小标题
\newcommand{\lithead}[1]{%
  \vspace{0.3em}%
  \noindent{\small\bfseries\color{titlecolor}#1}\par
  \vspace{0.1em}%
}

%% ── 环境 ──
\newenvironment{exampenv}[1]{%
  \vspace{0.5em}%
  \noindent{\color{supercolor!15}\rule{0.4em}{0pt}}\hspace{-0.4em}%
  \textcolor{supercolor}{\rule{0.4em}{0.9em}}\hspace{0.3em}%
  \textcolor{supercolor}{\textbf{【案例/例题】#1}}\par
  \noindent\hspace{0.7em}{\small\color{infocolor}分析：}\normalsize
}{%
  \par\vspace{0.3em}%
}

\newenvironment{pracenv}[1]{%
  \vspace{0.5em}%
  \noindent{\color{learncolor!15}\rule{0.4em}{0pt}}\hspace{-0.4em}%
  \textcolor{learncolor}{\rule{0.4em}{0.9em}}\hspace{0.3em}%
  \textcolor{learncolor}{\textbf{【练习/复盘】#1}}\par
  \noindent\hspace{0.7em}{\small\color{gray}解答：}\normalsize
}{%
  \par\vspace{0.3em}%
}

%% ══════════════════════════════════════════════════════════════════

%% ── 中文 ──
\usepackage{xeCJK}
\setCJKmainfont{SimSun}
\usepackage{indentfirst}

%% ── 数学/公式 ──
\usepackage{amsmath,amssymb,amsthm}
\usepackage{cancel}
\usepackage{physics}
\usepackage{mathrsfs}
\usepackage{bm}

%% ── 页面 ──
\usepackage[top=1.4cm,bottom=1.2cm,left=1.0cm,right=1.0cm,includehead]{geometry}
\usepackage{xcolor}
\usepackage{enumitem}
\usepackage{array,booktabs,multirow}
\usepackage{fancyhdr}
\usepackage{titlesec}
\usepackage{hyperref}
\usepackage{environ}
\usepackage{tikz}
\usetikzlibrary{decorations.pathmorphing,patterns,arrows.meta}
\hypersetup{colorlinks=true,linkcolor=blue!60!black,urlcolor=blue!60!black}

%% ── 颜色系统 ──
\definecolor{titlecolor}{HTML}{1a237e}
\definecolor{seccolor}{HTML}{004d40}
\definecolor{subseccolor}{HTML}{b71c1c}
\definecolor{emphcolor}{HTML}{bf360c}
\definecolor{infocolor}{HTML}{1565c0}
\definecolor{warncolor}{HTML}{e65100}
\definecolor{supercolor}{HTML}{880e4f}
\definecolor{examplecolor}{HTML}{1b5e20}
\definecolor{sectionbg}{HTML}{e8eaf6}
\definecolor{lightgray}{HTML}{f5f5f5}
\definecolor{hwcolor}{HTML}{795548}
\definecolor{learncolor}{HTML}{2e7d32}
\definecolor{papercolor}{HTML}{4a148c}   % 文献卡片紫

%% ── 标题格式 ──
\titleformat{\section}{\LARGE\bfseries\color{titlecolor}}{}{0em}{}[\vspace{-0.2em}\rule{\textwidth}{0.8pt}]
\titleformat{\subsection}{\large\bfseries\color{seccolor}}{}{0em}{}
\titleformat{\subsubsection}{\normalsize\bfseries\color{subseccolor}}{}{0em}{}
\titlespacing*{\section}{0pt}{1em}{0.5em}
\titlespacing*{\subsection}{0pt}{0.6em}{0.2em}
\titlespacing*{\subsubsection}{0pt}{0.4em}{0.1em}

%% ── 页眉（使用时替换标题） ──
\pagestyle{fancy}
\fancyhf{}
\fancyhead[L]{\small\color{gray}astro-nova · 文献笔记}
\fancyhead[R]{\small\color{gray}\today}
\fancyfoot[C]{\small\color{gray}\thepage}
\renewcommand{\headrulewidth}{0.4pt}
\setlength{\headheight}{14pt}

%% ── NovaForge 通用命令 ──
\newcommand{\key}[1]{\textcolor{emphcolor}{\textbf{#1}}}
\newcommand{\super}[1]{\textcolor{supercolor}{\textbf{#1}}}
\newcommand{\source}[1]{\textcolor{purple!60!black}{\small [#1]}}
\newcommand{\formula}[1]{\vspace{0.3em}\begin{center}\fcolorbox{titlecolor}{white}{\parbox{0.92\textwidth}{\centering\color{titlecolor}\small #1}}\end{center}\vspace{0.2em}}
\newcommand{\infobox}[1]{\vspace{0.3em}{\small\color{infocolor}\noindent$\blacktriangleright$ #1}\vspace{0.2em}}
\newcommand{\warning}[1]{\vspace{0.3em}{\small\color{warncolor}\noindent\textbf{!} #1}\vspace{0.2em}}
\newcommand{\sep}{\vspace{0.5em}\hrule\vspace{0.5em}}

\newcommand{\knowtitle}[1]{%
  \vspace{0.4em}%
  \noindent\colorbox{sectionbg}{\parbox{\dimexpr\textwidth-2\fboxsep\relax}{\small\bfseries\color{titlecolor}#1}}%
  \vspace{0.2em}%
}

%% ── 科研模式专有命令 ──

% 文献卡片
\newcommand{\paperinfo}[4]{%
  \vspace{0.5em}%
  \noindent
  \begin{tabular}{|p{2.2em}|p{\dimexpr\textwidth-2.2em-2\tabcolsep\relax}|}
    \hline
    \multicolumn{1}{|c|}{\textbf{\small\color{papercolor}标题}} & {\small #1}\\
    \hline
    \multicolumn{1}{|c|}{\textbf{\small\color{papercolor}作者}} & {\small #2}\\
    \hline
    \multicolumn{1}{|c|}{\textbf{\small\color{papercolor}刊源}} & {\small #3}\\
    \hline
    \multicolumn{1}{|c|}{\textbf{\small\color{papercolor}年份}} & {\small #4}\\
    \hline
  \end{tabular}\par
  \vspace{0.3em}%
}

% 科研小标题
\newcommand{\lithead}[1]{%
  \vspace{0.3em}%
  \noindent{\small\bfseries\color{titlecolor}#1}\par
  \vspace{0.1em}%
}

%% ── 环境 ──
\newenvironment{exampenv}[1]{%
  \vspace{0.5em}%
  \noindent{\color{supercolor!15}\rule{0.4em}{0pt}}\hspace{-0.4em}%
  \textcolor{supercolor}{\rule{0.4em}{0.9em}}\hspace{0.3em}%
  \textcolor{supercolor}{\textbf{【案例/例题】#1}}\par
  \noindent\hspace{0.7em}{\small\color{infocolor}分析：}\normalsize
}{%
  \par\vspace{0.3em}%
}

\newenvironment{pracenv}[1]{%
  \vspace{0.5em}%
  \noindent{\color{learncolor!15}\rule{0.4em}{0pt}}\hspace{-0.4em}%
  \textcolor{learncolor}{\rule{0.4em}{0.9em}}\hspace{0.3em}%
  \textcolor{learncolor}{\textbf{【练习/复盘】#1}}\par
  \noindent\hspace{0.7em}{\small\color{gray}解答：}\normalsize
}{%
  \par\vspace{0.3em}%
}
